\subsection{Injective cogenerator modules}

PROPOSITION 11. --- Let B be a ring, F a B-module, and P a (B, A)-bimodule. If F is an injective B-module and P a flat A-module, \(\mathrm{Hom}_{\mathbf{B}}(\mathbf{P},\mathbf{F})\) is an injective A-module.

Let \(u: \mathbf{M}' \to \mathbf{M}\) be an injective homomorphism of A-modules. We have a commutative diagram

\[
\begin{array}{c c c} \operatorname{Hom} _ {\mathbf {A}} (\mathbf {M}, \operatorname{Hom} _ {\mathbf {B}} (\mathbf {P}, \mathbf {F})) & \xrightarrow {\operatorname{Hom} _ {\mathbf {A}} (u , 1)} & \operatorname{Hom} _ {\mathbf {A}} (\mathbf {M} ^ {\prime}, \operatorname{Hom} _ {\mathbf {B}} (\mathbf {P}, \mathbf {F})) \\ \beta \Big \downarrow & & \beta^ {\prime} \Big \downarrow \\ \operatorname{Hom} _ {\mathbf {B}} (\mathbf {P} \otimes_ {\mathbf {A}} \mathbf {M}, \mathbf {F}) & \xrightarrow {\operatorname{Hom} (1 _ {\mathbf {P}} \otimes u , 1 _ {\mathbf {F}})} & \operatorname{Hom} _ {\mathbf {B}} (\mathbf {P} \otimes_ {\mathbf {A}} \mathbf {M} ^ {\prime}, \mathbf {F}) \end{array}
\]

where \(\beta\) and \(\beta'\) are the canonical isomorphisms of II, p.~74. Since P is flat over A, the homomorphism \(1_{\mathbf{P}} \otimes u: \mathbf{P} \otimes_{\mathbf{A}} \mathbf{M}' \to \mathbf{P} \otimes_{\mathbf{A}} \mathbf{M}\) is injective. Since F is injective, Hom ( \(1_{\mathbf{P}} \otimes u, 1_{\mathbf{F}}\) ) is surjective, hence also \(\operatorname{Hom}_{\mathbf{A}}(u, 1)\) , which proves that \(\operatorname{Hom}_{\mathbf{F}}(\mathbf{P}, \mathbf{F})\) is an injective A-module (X, p.~16, lemma 3).

DEFINITION 3. --- The A-module E is said to be a cogenerator if, for every A-module M and every nonzero element x of M, there exists an A-homomorphism u : M → E such that u(x) ≠ 0.

We say that the A-module L is a generator if, for every A-module M and every element x of M, there exists an A-homomorphism \(u : L \to M\) such that \(x \in u(L)\) . For example, the A-module \(A_{s}\) is a generator.

PROPOSITION 12. --- Let E be an injective A-module. For E to be a cogenerator, it is necessary and sufficient that \(\mathrm{Hom}_{\mathbf{A}}(\mathbf{S},\mathbf{E})\neq 0\) for every simple A-module S.

The condition is evidently necessary. Conversely, let M be an A-module and x a nonzero element of M; the submodule Ax of M has a simple quotient S (VIII, § 2, no. 1, prop. 3). If Hom \(_{A}\) (S, E) ≠ 0, then Hom \(_{A}\) (Ax, E) ≠ 0 and there exists a homomorphism f : Ax → E such that f(x) ≠ 0; since E is injective, f extends to a homomorphism u from M to E and we have u(x) = f(x) ≠ 0.

Example. --- The injective Z-module Q/Z (X, p.~18, example 2) is a cogenerator. Indeed, every simple Z-module is isomorphic to a module Z/pZ, \(p \neq 0\) , and \(\operatorname{Hom}_{\mathbf{Z}}(\mathbf{Z}/p\mathbf{Z}, \mathbf{Q}/\mathbf{Z})\) is nonzero (it contains, for example, the map induced by passing to quotients from the homomorphism \(x \mapsto x/p\) from Z to Q).

PROPOSITION 13. --- Let B be a ring, F an injective cogenerator B-module, P a (B, A)-bimodule. Suppose P is flat over A and such that P ⊗A S ≠ 0 for every simple A-module S (that is, faithfully flat over A in the sense of AC, I). Then the A-module Hom\_B (P, F) is a cogenerator and is injective.

Indeed, \(\mathrm{Hom}_{\mathbf{B}}(\mathbf{P},\mathbf{F})\) is injective by Prop. 11. On the other hand, for every simple A-module S, \(\mathrm{Hom}_{\mathbf{A}}(\mathbf{S},\mathrm{Hom}_{\mathbf{B}}(\mathbf{P},\mathbf{F}))\) is isomorphic to \(\mathrm{Hom}_{\mathbf{B}}(\mathbf{P}\otimes_{\mathbf{A}}\mathbf{S},\mathbf{F})\) therefore is nonzero since \(\mathbf{P}\otimes_{\mathbf{A}}\mathbf{S}\neq 0\) and the B-module F is a cogenerator; the A-module \(\mathrm{Hom}_{\mathbf{B}}(\mathbf{P},\mathbf{F})\) is therefore a cogenerator by Prop. 12.

Corollary 1. --- The A-module \(E_{A} = \operatorname{Hom}_{Z}(A, Q/Z)\) is injective and a cogenerator.

We apply Prop. 13 with \(\mathbf{B} = \mathbf{Z}\) , \(\mathbf{F} = \mathbf{Q} / \mathbf{Z}\) (example) and \(\mathbf{P} = \mathbf{A}_d\) .

For every A-module M, set

\[
\mathrm{I} ^ {0} (\mathbf {M}) = \mathrm{E} _ {\mathbf {A}} ^ {\mathrm{Hom} (\mathbf {M}, \mathbf {E} _ {\mathbf {A}})}
\]

and let \(e_{M}: M \to I^{0}(M)\) be the homomorphism which associates to \(m \in M\) the element

\[
\big (\varphi (m) \big) _ {\varphi \in \operatorname{Hom} (\mathbf {M}, \mathrm{E} _ {\mathbf {A}})} \in \mathrm{I} ^ {0} (\mathbf {M}).
\]

Then:

COROLLARY 2. --- The A-module I \(^{0}\) (M) is injective and the A-homomorphism \(e_{\mathbf{M}} : \mathbf{M} \to \mathbf{I}^{0}(\mathbf{M})\) is injective.

Indeed, I \(^{0}\) (M) is injective, since E \(_{A}\) is injective (X, p.~16, prop. 9), and on the other hand e \(_{M}\) is injective since E \(_{A}\) is a cogenerator.

COROLLARY 3. --- Every A-module is isomorphic to a submodule of an injective A-module.

COROLLARY 4. --- For the A-module E to be injective, it is necessary and sufficient that every injective A-homomorphism \(f: \mathbf{E} \to \mathbf{F}\) possess an A-linear retraction.

Suppose E is injective and let \(f: \mathrm{E} \to \mathrm{F}\) be an injective A-homomorphism. Then \(\operatorname{Hom}_{\mathbf{A}}(f, 1_{\mathbf{E}}): \operatorname{Hom}_{\mathbf{A}}(\mathrm{F}, \mathrm{E}) \to \operatorname{Hom}_{\mathbf{A}}(\mathrm{E}, \mathrm{E})\) is surjective; there therefore exists \(r \in \operatorname{Hom}_{\mathbf{A}}(\mathrm{F}, \mathrm{E})\) such that \(r \circ f = 1_{\mathbf{E}}\) and \(r\) is an A-linear retraction of \(f\) . Conversely, there exists an injective A-module I and an injective A-homomorphism \(f: \mathrm{E} \to \mathrm{I}\) (cor. 2); if \(f\) possesses an A-linear retraction, E is injective by prop. 9 of X, p.~16.

